% Part: proof-theory
% Chapter: sequent-calculus
% Section: invertibility

\documentclass[../../../include/open-logic-section]{subfiles}

\begin{document}

\olfileid[id]{pt}{seq}{inv}

\olsection{Keterbalikan Aturan}

\begin{defn}
Suatu aturan~$R$,
\[
\AxiomC{$S_1$}
\AxiomC{$\dots$}
\AxiomC{$S_n$}
\RightLabel{$R$}
\TrinaryInfC{$S$}
\DisplayProof
\]
disebut \emph{dapat dibalik} dalam suatu sistem jika, setiap kali
$\Proves S$, berlaku $\Proves[h_i] S_i$ untuk $i = 1$, \dots,~$n$.
Aturan tersebut disebut \emph{dapat dibalik dengan mempertahankan
!!{height}} (atau \emph{!!{hp}-dapat dibalik}) dalam suatu sistem jika,
setiap kali $\Proves_n S$, berlaku $\Proves[n] S_i$ untuk $i = 1$,
\dots,~$n$.
\end{defn}

\begin{lem}\ollabel{lem:G3c-invert}
Aturan-aturan~\Log{G3c} untuk $\lnot$, $\land$, $\lor$, dan $\lif$
dapat dibalik dengan mempertahankan !!{height}, yakni:
\begin{enumerate}
  \item \RightR{\lnot}: Jika $\Log{G3c} \Proves[n] \Gamma \Sequent
  \Delta, \lnot !A$, maka $\Log{G3c} \Proves[n] !A, \Gamma \Sequent
  \Delta$.
  \item \LeftR{\lnot}: Jika $\Log{G3c} \Proves[n] \Gamma, \lnot !A
  \Sequent \Delta$, maka $\Log{G3c} \Proves[n] \Gamma \Sequent \Delta,
  !A$.
  \item \RightR{\land}: Jika $\Log{G3c} \Proves[n] \Gamma \Sequent
  \Delta, !A \land !B$, maka $\Log{G3c} \Proves[n] \Gamma \Sequent
  \Delta, !A$ dan $\Log{G3c} \Proves[n] \Gamma \Sequent
  \Delta, !B$.
  \item \LeftR{\land}: Jika $\Log{G3c} \Proves[n] \Gamma, !A \land !B
  \Sequent \Delta$, maka $\Log{G3c} \Proves[n] !A, !B, \Gamma \Sequent
  \Delta$.
  \item \RightR{\lor}: Jika $\Log{G3c} \Proves[n] \Gamma \Sequent
  \Delta, !A \lor !B$, maka $\Log{G3c} \Proves[n] \Gamma \Sequent
  \Delta, !A, !B$.
  \item \LeftR{\lor}: Jika $\Log{G3c} \Proves[n] !A \lor !B, \Gamma \Sequent
  \Delta$, maka $\Log{G3c} \Proves[n] !A, \Gamma \Sequent
  \Delta$ dan $\Log{G3c} \Proves[n] !B, \Gamma \Sequent
  \Delta$.
  \item \RightR{\lif}: Jika $\Log{G3c} \Proves[n] \Gamma \Sequent
  \Delta, !A \lif !B$, maka $\Log{G3c} \Proves[n] !A, \Gamma \Sequent
  \Delta, !B$.
  \item \LeftR{\lif}: Jika $\Log{G3c} \Proves[n] !A \lif !B, \Gamma
  \Sequent \Delta$, maka $\Log{G3c} \Proves[n] \Gamma \Sequent \Delta,
  !A$ dan $\Log{G3c} \Proves[n] !B, \Gamma \Sequent \Delta$.
\end{enumerate}
\end{lem}

\begin{proof}
Kita harus menunjukkan bahwa, untuk setiap aturan~$R$ dari~\Log{G3c},
jika terdapat !!a{proof}~$\pi$ atas kesimpulan~$S$ dari~$R$, maka juga
terdapat !!{proof}s~$\pi_i$ atas premis-premis $S_i$ dari~$R$, dengan
$\pheight{\pi_i} \le \pheight{\pi}$. Tinjau aturan \LeftR{\land}:
andaikan terdapat !!a{proof}~$\pi$ dengan sekuen akhir berbentuk $!A
\land !B, \Gamma \Sequent \Delta$. Kita harus menunjukkan bahwa
terdapat !!a{proof} $\pi'$ atas $!A, !B, \Gamma \Sequent \Delta$
dengan $\pheight{\pi'} \le \pheight{\pi}$. Kita melakukannya dengan
induksi pada $n = \pheight{\pi}$.

Jika $n = 0$, $\pi$ hanya terdiri atas sebuah aksioma: $!A \land !B,
\Gamma \Sequent \Delta$. Dalam \Log{G3c}, aksioma berbentuk $!C,
\Gamma \Sequent \Delta, !C$ dengan $!C$ atomik. Dengan kata lain,
$!C$ tidak mungkin sama dengan $!A \land !B$. Jika $!A \land !B,
\Gamma \Sequent \Delta$ merupakan aksioma, suatu $!C$ atomik harus
merupakan !!a{element} dari $\Gamma$ sekaligus~$\Delta$. Namun demikian,
$!A, !B, \Gamma \Sequent \Delta$ juga merupakan aksioma, sehingga kita
telah memperoleh !!{proof}~$\pi'$ yang dicari (yang juga mempunyai
!!{height}~$0$).

Jika $n>0$, bukti tersebut berakhir dengan sebuah inferensi. Kita harus
membuktikan klaim untuk~$n$, tetapi boleh mengandaikan bahwa klaim itu
benar bagi setiap !!{proof} dengan ketinggian~$<n$, sekalipun konteksnya
berbeda. Dengan kata lain, untuk setiap $k<n$ serta setiap $\Pi$ dan
$\Lambda$, jika $!A \land !B, \Pi \Sequent \Lambda$ mempunyai
!!a{proof} dengan !!{height}~$k$, maka terdapat !!a{proof} dengan
!!{height}~$\le k$ atas $!A, !B, \Pi \Sequent \Lambda$.

Ada dua kasus: $!A \land !B$ di sebelah kiri bersifat prinsipal, atau
tidak. Jika bersifat prinsipal, inferensi terakhir adalah \LeftR{\land}
dan sub-!!{proof} langsung~$\pi_1$ berakhir pada $!A, !B, \Gamma
\Sequent \Delta$. Hasilnya berlaku dalam kasus ini karena
$\pheight{\pi_1} < \pheight{\pi}$.

Jika $!A \land !B$ tidak prinsipal, suatu !!{formula} lainlah yang
prinsipal, sedangkan $!A \land !B$ termasuk dalam konteks inferensi
terakhir. Hipotesis induksi lalu berlaku bagi sub-!!{proof}s langsung
dari aturan terakhir~$R$, dan menghasilkan !!a{proof}~$\pi_1'$ atau dua
!!{proof}s~$\pi_1'$ dan~$\pi_2'$ dengan !!{height} yang tidak lebih besar
daripada ketinggian sub-!!{proof}s langsung dari~$\pi$. Sekuen akhir
$\pi_1'$ dan $\pi_2'$ memuat $!A, !B$ dalam konteks di sebelah kiri,
menggantikan $!A \land !B$. Terapkan aturan~$R$ untuk memperoleh
bukti~$\pi'$. Sebagai contoh, andaikan $\pi$ berakhir dengan~\LeftR{\lif}:
\[
\AxiomC{}
\RightLabel{$\pi_1$}
\Deduce$!A \land !B, \Gamma' \fCenter \Delta, !C$
\AxiomC{}
\RightLabel{$\pi_2$}
\Deduce$!A \land !B, \Gamma', !D \fCenter \Delta$
\RightLabel{\LeftR{\lif}}
\BinaryInf$!A \land !B, \Gamma', !C \lif !D \fCenter \Delta$
\DisplayProof
\]
Di sini $!C \lif !D$ di sebelah kiri bersifat prinsipal, sedangkan
konteks di sebelah kiri ialah $!A \land !B, \Gamma'$ (yakni,
$\Gamma = \Gamma', !C \lif !D$). Hipotesis induksi menghasilkan
!!{proof}s $\pi_1'$ dan $\pi_2'$ masing-masing atas $!A, !B, \Gamma'
\Sequent \Delta, !C$ dan $!A, !B, \Gamma', !D \Sequent \Delta$.
Kedua sekuen ini kembali dapat menjadi premis aturan \LeftR{\lif}
untuk memperoleh~$\pi'$:
\[
\AxiomC{}
\RightLabel{$\pi_1'$}
\Deduce$!A, !B, \Gamma' \fCenter \Delta, !C$
\AxiomC{}
\RightLabel{$\pi_2'$}
\Deduce$!A, !B, \Gamma', !D \fCenter \Delta$
\RightLabel{\LeftR{\lif}}
\BinaryInf$!A, !B, \Gamma', !C \lif !D \fCenter \Delta$
\DisplayProof
\]
Kita mempunyai
\begin{align*}
\pheight{\pi} & = \max(\pheight{\pi_1}, \pheight{\pi_2}) + 1\\
\pheight{\pi'} & = \max(\pheight{\pi_1'}, \pheight{\pi_2'}) + 1
\intertext{menurut definisi, dan}
\pheight{\pi_1'} & \le \pheight{\pi_1}\\
\pheight{\pi_2'} & \le \pheight{\pi_2}
\intertext{menurut hipotesis induksi. Jadi,}
\pheight{\pi'} & \le \pheight{\pi}.
\end{align*}
\end{proof}

\begin{prob}
Berikan bukti \olref{lem:G3c-invert} untuk \RightR{\lif}.
\end{prob}

\begin{lem}\ollabel{lem:invert-quant}
Aturan \RightR{\lforall} dan \LeftR{\lexists} dapat dibalik dengan
mempertahankan !!{height} dalam arti berikut:
\begin{enumerate}
  \item Jika $\Log{G3c} \Proves[n] \Gamma \Sequent \Delta,
  \lforall[x][!A(x)]$, maka $\Log{G3c} \Proves[n] \Gamma \Sequent
  \Delta, !A(c)$, dan
  \item jika $\Log{G3c} \Proves[n] \lexists[x][!A(x)], \Gamma \Sequent
  \Delta$, maka $\Log{G3c} \Proves[n] !A(c), \Gamma \Sequent \Delta$,
\end{enumerate}
untuk setiap $c$ yang tidak muncul dalam $\Gamma$, $\Delta$, dan
$\lforall[x][!A(x)]$ atau $\lexists[x][!A(x)]$, secara berturut-turut.
\end{lem}

\begin{proof}
  Kita menunjukkan (1) dan menyerahkan (2) sebagai latihan. Argumennya
  mirip dengan bukti \olref{lem:G3c-invert}. Pada langkah induksi,
  kembali terdapat dua kasus. Kasus ketika $\lforall[x][!A(x)]$
  prinsipal segera selesai: premis inferensi $\RightR{\lforall}$ terakhir
  berbentuk sebagaimana diperlukan, yakni $\Gamma \Sequent \Delta,
  !A(a)$, dan sub-!!{proof} langsung~$\pi_1$ yang berakhir padanya
  memenuhi $\pheight{\pi_1} < \pheight{\pi}$. Selain itu, syarat
  variabel eigen memastikan bahwa $a$ tidak muncul dalam $\Gamma$,
  $\Delta$, ataupun $\lforall[x][!A(x)]$. Karena kita boleh mengandaikan
  $\pi_1$ reguler, $\Subst{\pi_1}{c}{a}$ merupakan !!a{proof} dengan
  !!{height} yang sama atas $\Gamma \Sequent \Delta, !A(c)$ menurut
  \olref[qua]{lem:subst-regular}.

  Dalam kasus kedua, $\lforall[x][!A(x)]$ tidak prinsipal dalam
  inferensi terakhir; suatu formula lainlah yang prinsipal. Kita kembali
  meninjau satu kasus sebagai contoh. Andaikan inferensi terakhir adalah
  ~\RightR{\land}. Suatu !!{formula} dalam $\Delta$ berbentuk~$!B \land
  !C$ dan bersifat prinsipal. Adapun !!{proof}~$\pi$ berakhir dengan
  \[
\AxiomC{}
\RightLabel{$\pi_1$}
\Deduce$\Gamma \fCenter \Delta', \lforall[x][!A(x)], !B$
\AxiomC{}
\RightLabel{$\pi_2$}
\Deduce$\Gamma \fCenter \Delta', \lforall[x][!A(x)], !C$
\RightLabel{\RightR{\land}}
\BinaryInf$\Gamma \fCenter \Delta', \lforall[x][!A(x)], !B \land !C$
\DisplayProof
\]
dengan $\Delta = \Delta', !B \land !C$. Ambil $c$ sebagai
!!a{constant} yang tidak muncul dalam $\Delta$; maka jelas bahwa konstanta itu
juga tidak muncul dalam $\Delta', !B$ maupun $\Delta', !C$. Jadi,
hipotesis induksi berlaku bagi $\pi_1$ dan~$\pi_2$; kita memperoleh
!!{proof}s~$\pi_1'$ dan $\pi_2'$ dengan $\pheight{\pi_1'} \le
\pheight{\pi_1}$ serta $\pheight{\pi_2'} \le \pheight{\pi_2}$.
Bentuk !!{proof}~$\pi'$ yang diperlukan atas $\Gamma \Sequent \Delta, !A(c)$
ialah:
\[
\AxiomC{}
\RightLabel{$\pi_1'$}
\Deduce$\Gamma \fCenter \Delta', !A(c), !B$
\AxiomC{}
\RightLabel{$\pi_2'$}
\Deduce$\Gamma \fCenter \Delta', !A(c), !C$
\RightLabel{\RightR{\land}}
\BinaryInf$\Gamma \fCenter \Delta', !A(c), !B \land !C$
\DisplayProof
\]
dengan $\pheight{\pi'} = \max(\pheight{\pi_1'}, \pheight{\pi_2'})+1
\le \pheight{\pi}$.
\end{proof}

\Olref{lem:G3c-invert} berperan penting dalam bukti teorema eliminasi
potongan. Namun, lema itu juga dapat dipakai untuk menunjukkan hasil berikut:

\begin{prop}\ollabel{prop:G3c-cont-adm}
Aturan kontraksi \LeftR{\Contraction} dan \RightR{\Contraction}
dari~\Log{G1c} bersifat admisibel dalam~\Log{G3c} dengan mempertahankan
!!{height}.
\end{prop}

\begin{proof}
Kita memberikan argumen untuk \RightR{\Contraction} dan
\LeftR{\Contraction} sekaligus. Andaikan $\pi$ merupakan !!a{proof} atas
$\Gamma \Sequent \Delta, !A, !A$ atau atas $!A, !A, \Gamma \Sequent
\Delta$ dalam~\Log{G3c}. Kita harus menunjukkan bahwa juga terdapat
!!a{proof}~$\pi'$ dalam \Log{G3c} atas $\Gamma \Sequent \Delta, !A$ atau $!A,
\Gamma \Sequent \Delta$, secara berturut-turut, dengan
$\pheight{\pi'} \le \pheight{\pi}$. Kita melakukannya dengan induksi
pada $n = \pheight{\pi}$.

Jika $n=0$, $\pi$ hanyalah sebuah aksioma. Namun, jika $\Gamma \Sequent
\Delta, !A, !A$ merupakan aksioma~\Log{G3c}, demikian pula $\Gamma
\Sequent \Delta, !A$: entah suatu $!C$ atomik merupakan !!a{element}
dari $\Gamma$ sekaligus~$\Delta$, atau $!A$ bersifat atomik dan merupakan
!!a{element} dari~$\Gamma$. Penalaran yang sama berlaku jika sekuen
akhirnya ialah $!A, !A, \Gamma \Sequent \Delta$.

Jika $n>0$, $\pi$ berakhir dengan suatu aturan~$R$. Jika $!A$ tidak
prinsipal dalam inferensi terakhir ini, kita dapat memperoleh
!!a{proof}~$\pi'$ seperti dalam bukti \olref{lem:G3c-invert}, yaitu dengan
menerapkan hipotesis induksi pada sub-!!{proof}s langsung dari~$\pi$ dan
menerapkan aturan~$R$ yang sama pada satu atau dua !!{proof} yang dihasilkan.
Hipotesis induksinya ialah bahwa jika terdapat !!a{proof} dengan
!!{height}~$k<n$ atas $\Pi \Sequent \Lambda, !D, !D$ atau atas $!D, !D,
\Pi \Sequent \Lambda$, maka terdapat !!a{proof} dengan !!{height}~$\le k$
atas $\Pi \Sequent \Lambda, !D$ atau $!D, \Pi \Sequent \Lambda$, secara
berturut-turut. (Perhatikan bahwa hipotesis induksi tetap berlaku jika
konteks atau !!{formula} yang dikontraksi berbeda dari $\Gamma$, $\Delta$,
atau~$!A$!)

Sebagai contoh, andaikan $R$ adalah $\RightR{\land}$ dan $\pi$~berakhir
dengan
\[
\AxiomC{}
\RightLabel{$\pi_1$}
\Deduce$\Gamma \fCenter \Delta', !B, !A, !A$
\AxiomC{}
\RightLabel{$\pi_2$}
\Deduce$\Gamma \fCenter \Delta', !C, !A, !A$
\RightLabel{\RightR{\land}}
\BinaryInf$\Gamma \fCenter \Delta', !B \land !C, !A, !A$
\DisplayProof
\]
dengan $\Delta = \Delta', !B \land !C$. Hipotesis induksi menghasilkan
!!{proof}s $\pi_1'$ dan $\pi_2'$ atas $\Gamma \fCenter \Delta', !B, !A$
dan $\Gamma \fCenter \Delta', !C, !A$, secara berturut-turut, dengan
$\pheight{\pi_1'} \le \pheight{\pi_1}$ dan $\pheight{\pi_2'} \le
\pheight{\pi_2}$. Adapun !!{proof}~$\pi'$ ialah:
\[
\AxiomC{}
\RightLabel{$\pi_1'$}
\Deduce$\Gamma \fCenter \Delta', !B, !A$
\AxiomC{}
\RightLabel{$\pi_2'$}
\Deduce$\Gamma \fCenter \Delta', !C, !A$
\RightLabel{\RightR{\land}}
\BinaryInf$\Gamma \fCenter \Delta', !B \land !C, !A$
\DisplayProof
\]

Jika $!A$ prinsipal dalam inferensi terakhir, kita menerapkan lema
keterbalikan (\olref{lem:G3c-invert}) pada sub-!!{proof}s langsung
dari~$\pi$, kemudian hipotesis induksi, lalu aturan~$R$ sekali lagi.
Sebagai contoh, andaikan $\pi$ membuktikan $\Gamma \Sequent \Delta,
!A, !A$, dengan $!A \ident (!B \land !C)$ dan formula itu prinsipal dalam
inferensi terakhir. Maka $\pi$ berbentuk sebagai berikut:
\[
\AxiomC{}
\RightLabel{$\pi_1$}
\Deduce$\Gamma \fCenter \Delta, !B \land !C, !B$
\AxiomC{}
\RightLabel{$\pi_2$}
\Deduce$\Gamma \fCenter \Delta, !B \land !C, !C$
\RightLabel{\RightR{\land}}
\BinaryInf$\Gamma \fCenter \Delta, !B \land !C, !B \land !C$
\DisplayProof
\]
Terapkan \olref{lem:G3c-invert} pada $\pi_1$ untuk memperoleh
!!a{proof}~$\pi_1'$ atas $\Gamma \Sequent \Delta, !B, !B$. (Penerapan
itu juga menghasilkan !!a{proof} atas $\Gamma \Sequent \Delta, !C,
!B$, tetapi kita tidak memerlukannya.) Demikian pula, dengan menerapkan
\olref{lem:G3c-invert} pada~$\pi_2$, kita memperoleh !!a{proof}~$\pi_2'$
atas $\Gamma \Sequent \Delta, !C, !C$. Karena pembalikan
\RightR{\land} mempertahankan !!{height}, $\pheight{\pi_1'} \le
\pheight{\pi_1} < \pheight{\pi}$ dan $\pheight{\pi_2'} \le
\pheight{\pi_2} < \pheight{\pi}$. Karena itu, hipotesis induksi berlaku
bagi $\pi_1'$ dan $\pi_2'$: terdapat !!{proof}s~$\pi_1''$ dan $\pi_2''$
masing-masing atas $\Gamma \Sequent \Delta, !B$ dan $\Gamma \Sequent
\Delta, !C$, dengan $\pheight{\pi_1''} \le \pheight{\pi_1'} \le
\pheight{\pi_1}$ dan $\pheight{\pi_2''} \le \pheight{\pi_2'} \le
\pheight{\pi_2}$. Adapun !!{proof}~$\pi'$ kemudian ialah:
\[
\AxiomC{}
\RightLabel{$\pi_1''$}
\Deduce$\Gamma \fCenter \Delta, !B$
\AxiomC{}
\RightLabel{$\pi_2''$}
\Deduce$\Gamma \fCenter \Delta, !C$
\RightLabel{\RightR{\land}}
\BinaryInf$\Gamma \fCenter \Delta, !B \land !C$
\DisplayProof
\]
dengan $\pheight{\pi'} \le \pheight{\pi}$.

Kita harus memberikan perhatian khusus kepada aturan kuantor dalam
kasus-kasus ketika !!{formula} tersebut prinsipal.

Andaikan $!A \ident \lforall[x][!B(x)]$ prinsipal di sebelah kanan. Maka
$\pi$ berakhir dengan
\[
\AxiomC{}
\RightLabel{$\pi_1$}
\Deduce$\Gamma \fCenter \Delta, \lforall[x][!B(x)], !B(a)$
\RightLabel{\RightR{\lforall}}
\UnaryInf$\Gamma \fCenter \Delta, \lforall[x][!B(x)], \lforall[x][!B(x)]$
\DisplayProof
\]
Menurut \olref{lem:invert-quant}, terdapat !!a{proof}~$\pi_1'$ dengan
!!{height} $\le \pheight{\pi_1}$ atas $\Gamma \fCenter \Delta, !B(c),
!B(a)$ untuk setiap~$c$ yang tidak muncul dalam $\Gamma \fCenter \Delta,
\lforall[x][!B(x)], !B(a)$. Kita boleh mengandaikan bahwa $\pi_1'$
reguler; maka, menurut \olref[qua]{lem:subst-regular}, !!{proof}
$\Subst{\pi_1'}{a}{c}$ merupakan !!a{proof} atas $\Gamma \fCenter
\Delta, !B(a), !B(a)$ yang juga memiliki !!{height}~$\le
\pheight{\pi_1}$. Kini hipotesis induksi berlaku dan memberi kita
!!a{proof}~$\pi_1''$ atas $\Gamma \fCenter \Delta, !B(a)$. Terapkan
aturan \RightR{\lforall} untuk memperoleh~$\pi'$:
\[
\AxiomC{}
\RightLabel{$\pi_1''$}
\Deduce$\Gamma \fCenter \Delta, !B(a)$
\RightLabel{\RightR{\lforall}}
\UnaryInf$\Gamma \fCenter \Delta, \lforall[x][!B(x)]$
\DisplayProof
\]
dengan $\pheight{\pi'} \le \pheight{\pi}$.

Sekarang andaikan $!A \ident \lexists[x][!B(x)]$ dan formula itu prinsipal
di sebelah kanan. Maka $\pi$ berbentuk sebagai berikut:
\[
\AxiomC{}
\RightLabel{$\pi_1$}
\Deduce$\Gamma \fCenter \Delta, \lexists[x][!B(x)], \lexists[x][!B(x)], !B(t)$
\RightLabel{\RightR{\lexists}}
\UnaryInf$\Gamma \fCenter \Delta, \lexists[x][!B(x)], \lexists[x][!B(x)]$
\DisplayProof
\]
Hipotesis induksi berlaku bagi~$\pi_1$, sehingga terdapat
!!a{proof}~$\pi_1'$ atas $\Gamma \Sequent \Delta, \lexists[x][!B(x)],
!B(t)$. (Perhatikan bahwa lema keterbalikan tidak diperlukan di sini!)
Adapun !!{proof}~$\pi'$ kemudian ialah:
\[
\AxiomC{}
\RightLabel{$\pi_1'$}
\Deduce$\Gamma \fCenter \Delta, \lexists[x][!B(x)], !B(t)$
\RightLabel{\RightR{\lexists}}
\UnaryInf$\Gamma \fCenter \Delta, \lexists[x][!B(x)]$
\DisplayProof
\]
dengan $\pheight{\pi'} \le \pheight{\pi}$.
\end{proof}

\begin{prob}
Uraikan garis besar bukti \olref{prop:G3c-cont-adm} untuk
\LeftR{\Contraction}. Jelaskan kasus-kasus ketika $!A \ident (!B \land
!C)$ dan $!A \ident (!B \lif !C)$.
\end{prob}

\begin{prob}
Uraikan garis besar bukti \olref{prop:G3c-cont-adm} untuk
\LeftR{\Contraction} bagi kasus-kasus kuantor yang tersisa, yakni ketika
$!A \ident \lforall[x][!B(x)]$ dan $!A \ident \lexists[x][!B(x)]$.
\end{prob}


\end{document}
